\chapter{Plotting Polar Coordinates and Transforming them}
\lecture{2}{16 Apr. 14:16}{Handling Polar Coordinates}
\label{chap:plotting-polar-coordinates-and-transforming-them}

\section{Graph of Polar Coordinate System}
\label{sec:graph-polar-coord-sys}
\begin{figure}[H]
  \centering
  \begin{polargrid}[4][4][12][4.5][1]
  \end{polargrid}
  \caption{Polar Coordinate System}
  \label{fig:polar-coord-sys-graph}
\end{figure}

\begin{exercise}
  Plot:
  \begin{enumerate}
    \item \((-1,\frac{3\pi}{4})\)
    \item \((2,-\frac{3\pi}{2})\)
    \item \((-4,-\frac{5\pi}{6})\)
  \end{enumerate}
\end{exercise}
\begin{answer}
  \begin{polargrid}[4][4][24][4.5][1]
    \polarpoint{-1}{3*pi/4}[red][$(-1, \frac{3\pi}{4})$]
    \polarpoint{2}{-3*pi/2}[red][$(2, -\frac{3\pi}{2})$]
    \polarpoint{-4}{-5*pi/6}[red][$(-4, -\frac{5\pi}{6})$]
  \end{polargrid}
\end{answer}

\section{Deceptive Coordinates}
\label{sec:polar-deceptive-coord}
\begin{problem}[Same point in different form give different analysis]
  Does the point \((2,\frac{\pi}{2})\) lie on the curve \(r=2\cos 2\theta\)?\\
  How do you check it?
  \begin{answer}
    By substituting the value of \(r\) and \(\theta\) in the above equation.
    The answer is \textbf{NO}!
  \end{answer}
  But take the same point but written as: \((-2,\frac{\pi}{2})\). Now does it
  lie on the curve \(r=2\cos 2\theta\)? \textbf{YES}!
\end{problem}

\begin{note}[Solution to Deceptive Coordinates]
  \textbf{Use Cartesian Coordinates}!!\\
  To solve this duality of polar coordinates, simply change the polar
  coordinates to Cartesian Coordinate
  \autoref{sec:relating-cartesian-and-polar-coordinates}.
\end{note}


\section{Relating Cartesian and Polar Coordinates}
\label{sec:relating-cartesian-and-polar-coordinates}
Polar and Cartesian Coordinates relate by the following formulae:
\begin{enumerate}
  \item \(x^2+y^2 = r^2\) 
  \item \(\tan \theta = \frac{y}{x}\) 
  \item \(x = r\cos\theta\) 
  \item \(y = r\sin\theta\)
\end{enumerate}
\subsection*{Examples of Coordinate Conversion}
\label{ssec:examples-of-coordinate-conversion}

\begin{eg}[Cartesian to Polar]
  Convert the Cartesian coordinates $(3, 4)$ to polar coordinates $(r, \theta)$.

  \textbf{Solution:}
  \begin{itemize}
    \item Calculate radius: $r = \sqrt{3^2 + 4^2} = \sqrt{9 + 16} = \sqrt{25} = 5$
    \item Calculate angle: $\theta = \arctan\left(\frac{4}{3}\right) \approx 53.13^\circ$ or $0.93$ radians
  \end{itemize}
  Result: $(5, 53.13^\circ)$ or $(5, 0.93\ \text{rad})$
\end{eg}

\begin{eg}[Polar to Cartesian]
  Convert the polar coordinates $(2, \pi/3)$ to Cartesian coordinates $(x, y)$.

  \textbf{Solution:}
  \begin{itemize}
    \item Calculate $x$: $x = 2\cos\left(\frac{\pi}{3}\right) = 2 \times \frac{1}{2} = 1$
    \item Calculate $y$: $y = 2\sin\left(\frac{\pi}{3}\right) = 2 \times \frac{\sqrt{3}}{2} = \sqrt{3}$
  \end{itemize}
  Result: $(1, \sqrt{3})$
\end{eg}

\begin{remark}
  Remember these key points:
  \begin{itemize}
    \item For Cartesian \(\to \) Polar, always consider the \textbf{quadrant} when calculating $\theta$
    \item For Polar \(\to \) Cartesian, exact values are preferred when angles are special fractions of $\pi$
  \end{itemize}
\end{remark}

